\sscp{Cuscus \it{(Phalangeridae)}}{C-06-01}
Terms for cuscus (\it{Phalangeridae}) in \tsc{Seram-Tanimbar-Bomberai}
languages are given in here.
\ili{As} discussed in \srf{14-04-03-02},
we propose that \ili{Fordata}
and \ili{Kei} share with \ili{Banda}
and \ili{Aru} a local irregular development of penultimate *a > **ə;
\mbox{**mantəR} > **məntəR.

%\noindent{\wg\st{2pt}\begin{tabularx}{\linewidth}
%{lllll>{\raggedright\arraybackslash}X}
%\lsptoprule
%%Language	&	ISO	&	Term	&	Gloss	&	Etymon	&	Source	\\	\midrule
%%\endfirsthead
%Language	&	ISO	&	Term	&	Gloss	&	Etymon	&	Source	\\	\midrule
%\endhead
%
%\lspbottomrule
%\end{tabularx}}

\noindent{\wg\st{2pt}\begin{xltabular}{\linewidth}
{lllll>{\raggedright\arraybackslash}X}
\lsptoprule
%Language	&	ISO	&	Term	&	Gloss	&	Etymon	&	Source	\\	\midrule
%\endfirsthead
Language	&	ISO	&	Term	&	Gloss	&	Etymon	&	Source	\\	\midrule
\endhead
\mc{6}{c}{\tsc{East Seram}}		\\	\midrule	\hhline{~}
\ili{Masiwang}	&	bnf	&	mani	&	cuscus	&	*manuk	&	\cite[358]{LeCocqdArmandville1901}	\\	\midrule	\hhline{~}
\mc{6}{c}{\tsc{Eastern Islands}}		\\	\midrule	\hhline{~}
\ili{Bati}	&	bvt	&	manuk faniki	&	cuscus	&	*manuk	&	\cite[196]{Ellen1993}	\\
\ili{Geser}	&	ges	&	kidor	&	cuscus	&	**kindoR	&	\cite[196]{Ellen1993}	\\
\ili{Gorom}	&	ges	&	iˈdoʔ;	&	cuscus	&	**kindoR	&	\cite{Visser2019a};	\\	\hhline{~}
	&		&	‹i'do›	&	cuscus	&	**kindoR	&	\cite[51]{Kakerissa1986}	\\
\ili{Watubela}	&	wah	&	kadola	&	cuscus	&	**kVndoR(a)	&	\cite[241]{Blust1982} from Collins	\\
\ili{Kesui}	&	wah	&	udora	&	cuscus	&	**kVndoR(a)	&	\cite[241]{Blust1982} from Collins	\\
\ili{Kowiai}	&	kwh	&	indor	&	cuscus	&	**kindoR	&	\cite{WalkerWalker1991}	\\	\hhline{~}
\mc{6}{c}{\tsc{\ili{Teor}-Kur}}		\\	\midrule	\hhline{~}
\ili{Teor}	&	tev	&	medar	&	cuscus	&	**məntəR	&	Craig Marshall pers. comm. September 2025	\\	\midrule	\hhline{~}
\mc{6}{c}{\tsc{Tanimbar-Bomberai}}		\\	\midrule	\hhline{~}
\ili{Yamdena}\footnotemark 	&	jmd	&	manr-e	&	(white) &	**məntəR	&	\hp{\,}	\\	\hhline{~}
					&			&					&	cuscus	&						&	\multirow{-2}{=}{\raggedright \cite{MettlerMettler1996}, \cite[65]{Drabbe1932b}}	\\
\ili{Fordata}	&	frd	&	medar	&	cuscus	&	**məntəR	&	\cite{MarshallMarshall-Grimes1990}	\\
\ili{Fordata}, 	&	frd	&	kandruə	&	rat	&	**kVndoR(a)/	&	\hp{\,}	\\	\hhline{~}
\ili{Sera}\footnotemark			&			&					&			&	**kazupay	&	\multirow{-2}{=}{\raggedright \cite{MarshallMarshall-Grimes1990}}		\\
\ili{Kei} Kecil	&	kei	&	medar;	&	cuscus;	&	**məntəR	&	\cite{Travis2019}	\\	\hhline{~}
					&			&	leman		&	tan-coloured &	\it{misc.}	&		\\	\hhline{~}
					&			&					&	cuscus	&			&		\\
\ili{Kei} Besar	&	kei	&	mendar	&	cuscus	&	**məntəR	&	\cite[196]{Ellen1993}	\\
\ili{Uruangnirin}	&	urn	&	pora	&	cuscus	&	\it{misc.}	&	\cite{Visser2023}	\\
\ili{Onin}	&	oni	&	ket	&	cuscus	&	\it{misc.}	&	\cite{Donohue2010}	\\
\ili{Sekar}	&	skz	&	ket	&	cuscus	&	\it{misc.}	&	\cite{Donohue2010}	\\
\ili{Arguni}	&	agf	&	ʔara waked	&	cuscus	&	**waked	&	\cite{Narfafan2011}	\\
\ili{Erokwanas}	&	erw	&	matafeis	&	cuscus	&	**mantəR	&	\cite{Donohue2010}	\\
\ili{Irarutu}	&	irh	&	ʤemug	&	cuscus	&	\it{misc.}	&	\cite[224]{Jackson2014}	\\
\ili{Kuri}	&	nbn	&	auˈɛs	&	opossum	&	\it{misc.}	&	SV \citeyear[248]{SmitsVoorhoeve1992b}	\\
\lspbottomrule
\end{xltabular}}
\footnotetext[3]{
		\citet{MettlerMettler1996} give ‹manre› as ‘kuskus,
		putih’ = ‘cuscus, white’.
		\citet[65]{Drabbe1932b} gives ‹mande› as ‘boombeer’ = ‘tree bear’.
		The medial consonant in these forms is phonetically
		[ndr] analysed by \citet{MettlerMettler1997} as a prenasalised
		trill /nr/ realised as [nr], [ndr], or [nd].
		This form is probably analysable as containing the nominal suffix
		\it{-e} which replaces word-final a{\#} (see third paragraph
		of \srf{11-01-02}).}
\footnotetext{
		From the sound correspondences \ili{Sera} \it{kandruə} ‘rat’ appears to be a reflex of
		*kandoRa with metathesis, but could also be from PCEMP *kazupay ‘rat’.}